\documentclass[conference]{IEEEtran}
\IEEEoverridecommandlockouts

\usepackage{cite}
\usepackage{amsmath,amssymb,amsfonts}
\usepackage{algorithmic}
\usepackage{graphicx}
\usepackage{textcomp}
\usepackage{xcolor}
\usepackage{booktabs}
\usepackage{hyperref}
\usepackage{url}

\def\BibTeX{{\rm B\kern-.05em{\sc i\kern-.025em b}\kern-.08em
    T\kern-.1667em\lower.7ex\hbox{E}\kern-.125emX}}

\begin{document}

\title{Learning from Demonstration on ALE/MsPacman-v5:\\
A Three-Algorithm Comparison of DQN, Behavioral Cloning, and GAIL}

\author{
  \IEEEauthorblockN{Juan Diego Lozano Luna}
  \IEEEauthorblockA{\textit{20231020040}}
  \and
  \IEEEauthorblockN{Daniel Esteban Camacho Ospina}
  \IEEEauthorblockA{\textit{20231020046}}
  \and
  \IEEEauthorblockN{Juan Sebasti\'{a}n Rodr\'{i}guez Carre\~{n}o}
  \IEEEauthorblockA{\textit{20231020107}}
  \\[6pt]
  \IEEEauthorblockA{
    \textit{Machine Learning} \\
    \textit{Systems Engineering Program, Faculty of Engineering} \\
    \textit{Universidad Distrital Francisco Jos\'{e} de Caldas} \\
    Bogot\'{a}, Colombia
  }
}

\maketitle

\begin{abstract}
This paper presents the third installment of a sequential algorithm comparison series on
the Arcade Learning Environment (ALE) game \textit{Ms.\ Pac-Man} (\texttt{ALE/MsPacman-v5}).
Building upon a Deep Q-Network (DQN) agent from Challenge~1, we implement and evaluate
two imitation-learning approaches: Behavioral Cloning (BC) and Generative Adversarial
Imitation Learning (GAIL) with Proximal Policy Optimization (PPO) as the inner
reinforcement learning algorithm.
We investigate the central hypothesis that, on a dense-reward game such as Ms.\ Pac-Man,
the adversarial proxy reward of GAIL may underperform direct reward maximization, whereas
BC may prove surprisingly competitive by directly imitating high-quality demonstrations.
Experimental results confirm the hypothesis in a striking way: BC achieves a mean episode
return of $866.0 \pm 280.2$, outperforming both the baseline DQN ($85.6 \pm 15.6$) and
GAIL ($70.0 \pm 0.0$ under both demonstration sizes), which collapses to a degenerate
policy due to discriminator saturation.
A demonstration-size ablation (5\,000 vs.\ 50\,000 state-action pairs) reveals that GAIL
performance is insensitive to dataset size in the collapse regime.
These findings indicate that (i)~dense-reward games expose the fragility of adversarial
imitation learning under limited compute, (ii)~BC is a robust and computationally efficient
baseline when high-quality demonstrations are available, and (iii)~discriminator health
monitoring is critical for diagnosing GAIL training failure.
\end{abstract}

\begin{IEEEkeywords}
Generative Adversarial Imitation Learning, Behavioral Cloning, Deep Q-Network,
Proximal Policy Optimization, Atari, Ms.\ Pac-Man, Imitation Learning,
Discriminator Collapse
\end{IEEEkeywords}

% ──────────────────────────────────────────────────────────────────────────────
\section{Introduction}
\label{sec:intro}

The Arcade Learning Environment (ALE)~\cite{bellemare2013ale} provides a standardized
benchmark for evaluating reinforcement learning (RL) algorithms across dozens of video
games.
Over the course of this three-challenge series, we have studied three distinct learning
paradigms on the same environment, \texttt{ALE/MsPacman-v5}:

\begin{enumerate}
  \item \textbf{Challenge~1 — DQN}~\cite{mnih2015dqn}: value-based off-policy RL with
        experience replay.
  \item \textbf{Challenge~3 — PPO}~\cite{schulman2017ppo}: on-policy actor-critic with
        clipped surrogate objective.
  \item \textbf{Challenge~4 — GAIL}~\cite{ho2016gail}: imitation learning via adversarial
        training, augmented with a BC baseline.
\end{enumerate}

GAIL~\cite{ho2016gail} frames imitation as a two-player minimax game between a policy
(the generator) and a discriminator that learns to distinguish expert transitions from
agent-generated ones.
The adversarial reward replaces the true environment reward, making GAIL particularly
appealing for sparse-reward settings where standard RL struggles.
However, this reward substitution carries risk on dense-reward games: the adversarial proxy
may be weaker than the dense environment signal that PPO and DQN can exploit directly.

The central research question of this challenge is:
\textit{Does an agent that learns by imitating demonstrations converge faster, transfer
better, or reach higher final performance than pure RL agents (DQN, PPO) on Ms.\ Pac-Man?
Under which conditions does the demonstration source matter?}

Our specific group hypothesis (Group~7) states that Ms.\ Pac-Man's relatively dense reward
structure places GAIL at a disadvantage, and that BC alone --- which directly imitates
good movement patterns without any RL overhead --- may be competitive with or superior to
full GAIL on this game.
The experimental results reported here provide strong empirical support for this hypothesis.

The rest of the paper is organized as follows.
Section~\ref{sec:background} reviews the theoretical foundations of BC and GAIL.
Section~\ref{sec:environment} describes the environment and preprocessing pipeline.
Section~\ref{sec:method} details the implementation of each component.
Section~\ref{sec:experiments} presents the experimental protocol.
Section~\ref{sec:results} reports and analyzes the results.
Section~\ref{sec:discussion} addresses the four required analysis questions.
Section~\ref{sec:conclusion} concludes.

% ──────────────────────────────────────────────────────────────────────────────
\section{Background and Related Work}
\label{sec:background}

\subsection{Behavioral Cloning}

Behavioral Cloning (BC) treats imitation learning as supervised classification.
Given a demonstration dataset $\mathcal{D} = \{(s_i, a_i)\}_{i=1}^{N}$, BC minimises the
negative log-likelihood of the expert actions under the learnt policy $\pi_\theta$:

\begin{equation}
  \mathcal{L}_{\text{BC}}(\theta) = -\frac{1}{N} \sum_{i=1}^{N} \log \pi_\theta(a_i \mid s_i).
  \label{eq:bc}
\end{equation}

BC is computationally efficient and easy to implement.
Its principal weakness is \emph{distributional shift}~\cite{ross2011dagger}: at test time,
the policy visits states not covered by the demonstrations, and compounding errors drive it
to out-of-distribution regions.
DAgger~\cite{ross2011dagger} addresses this with iterative dataset aggregation, but at the
cost of additional environment interaction.

\subsection{Generative Adversarial Imitation Learning}

GAIL~\cite{ho2016gail} frames imitation as a two-player game between a policy $\pi_\theta$
(the generator) and a discriminator $D_\phi$.
The discriminator is trained to separate expert transitions from policy transitions using
binary cross-entropy:

\begin{equation}
  \max_{D_\phi}
  \mathbb{E}_{(s,a)\sim\mathcal{D}}[\log D_\phi(s,a)]
  + \mathbb{E}_{(s,a)\sim\pi_\theta}[\log(1 - D_\phi(s,a))].
  \label{eq:disc}
\end{equation}

The policy is trained with any RL algorithm using the \emph{adversarial reward}:

\begin{equation}
  r_{\text{adv}}(s,a) = \log D_\phi(s,a),
  \label{eq:adv_reward}
\end{equation}

which encourages the policy to produce trajectories that the discriminator cannot distinguish
from expert demonstrations.
The complete GAIL objective is:

\begin{equation}
  \min_{\pi_\theta} \max_{D_\phi}
  \mathbb{E}_{(s,a)\sim\pi_\theta}[\log D_\phi(s,a)]
  + \mathbb{E}_{(s,a)\sim\mathcal{D}}[\log(1-D_\phi(s,a))]
  - \lambda H(\pi_\theta),
  \label{eq:gail}
\end{equation}

where $H(\pi_\theta)$ is the policy entropy bonus (same as PPO's entropy coefficient).
Training alternates between a discriminator step (BCE on expert and agent mini-batches)
and a policy step (PPO update using $r_{\text{adv}}$).
The true environment reward is completely replaced during GAIL training and used only for
evaluation.

\subsection{Related Work}

Ho and Ermon~\cite{ho2016gail} first demonstrated GAIL on MuJoCo locomotion tasks with
low-dimensional state spaces.
Subsequent work has extended GAIL to image-based environments~\cite{torabi2018bco},
addressed reward shaping~\cite{fu2018airl}, and studied discriminator stability.
Ross et al.~\cite{ross2011dagger} showed that BC suffers from compounding errors but
proposed DAgger as a fix; our results suggest that in practice, with sufficient
demonstration coverage, BC can be remarkably effective.
Mnih et al.~\cite{mnih2015dqn} established the DQN baseline; Schulman et al.~\cite{schulman2017ppo}
introduced PPO as a stable on-policy alternative.

% ──────────────────────────────────────────────────────────────────────────────
\section{Environment and Preprocessing}
\label{sec:environment}

\subsection{ALE/MsPacman-v5}

Ms.\ Pac-Man is an Atari 2600 game with a 9-action discrete space.
The player navigates a maze eating pellets (+10 each) and power pellets (+50 each),
chasing/fleeing ghosts, and progressing through increasingly difficult maze layouts.
Unlike many ALE games, Ms.\ Pac-Man provides \emph{relatively dense rewards}: pellets
are distributed throughout the maze, so a random policy can still collect some reward.
This density is the key property motivating our hypothesis.

\subsection{Observation Preprocessing}

Preprocessing is identical across all three challenges to ensure comparability:
\begin{itemize}
  \item Noop reset (up to 30 no-ops at episode start).
  \item Frame skip of 4 (repeat action for 4 frames, return max-pooled).
  \item Grayscale conversion and resize to $84 \times 84$ pixels.
  \item Frame stack of 4 consecutive frames: observation shape $(4, 84, 84)$, \texttt{uint8}.
  \item Pixel values kept in $[0, 255]$; normalisation to $[0,1]$ is performed inside the
        network's forward pass.
\end{itemize}

The Gymnasium wrappers used are \texttt{AtariPreprocessing} and
\texttt{FrameStackObservation}, matching the configuration from Challenges 1 and 3.

% ──────────────────────────────────────────────────────────────────────────────
\section{Methodology}
\label{sec:method}

\subsection{Policy and Value Network (AtariActorCritic)}

Both the BC policy and the GAIL policy share the same Nature DQN convolutional
architecture~\cite{mnih2015dqn}:

\begin{itemize}
  \item \textbf{Conv1}: 32 filters, $8\times8$ kernel, stride 4, ReLU.
  \item \textbf{Conv2}: 64 filters, $4\times4$ kernel, stride 2, ReLU.
  \item \textbf{Conv3}: 64 filters, $3\times3$ kernel, stride 1, ReLU, Flatten.
  \item \textbf{FC}: Linear($64\!\times\!7\!\times\!7 \to 512$), ReLU.
  \item \textbf{Actor head}: Linear($512 \to 9$) — action logits.
  \item \textbf{Critic head}: Linear($512 \to 1$) — state value.
\end{itemize}

Total parameters: approximately 1.7M.

\subsection{GAIL Discriminator (GAILDiscriminator)}

The discriminator shares the same CNN backbone as the policy and outputs
$P(\text{expert} \mid s) \in (0,1)$:

\begin{itemize}
  \item CNN (same as above): 3136-dimensional feature vector.
  \item \textbf{FC1}: Linear($3136 \to 512$), Tanh.
  \item \textbf{FC2}: Linear($512 \to 1$), Sigmoid.
\end{itemize}

A state-only variant (\texttt{use\_action=False}) is used as the primary configuration.
The state-action variant (appending a one-hot action encoding) is available for future
ablation studies.

\subsection{Demonstration Collection}

Demonstrations are collected by rolling out the best DQN checkpoint from Challenge~1
(\texttt{exp\_02\_lr\_high}) deterministically (greedy action selection) and recording
$(s_t, a_t)$ pairs.
Two dataset sizes are collected for the ablation study:
\begin{itemize}
  \item \textbf{Primary}: 50\,000 steps (296 complete episodes).
  \item \textbf{Small} (ablation): 5\,000 steps (29 complete episodes).
\end{itemize}

All observations are stored channel-first ($(4, 84, 84)$, \texttt{uint8}) and compressed
with \texttt{np.savez\_compressed}.
Metadata (checkpoint path, step count, mean/std return) is written to a plain-text info
file as required by the deliverable specification.

\subsection{Behavioral Cloning Training}

The BC model minimises cross-entropy loss (Eq.~\ref{eq:bc}) over the primary 50k
demonstration dataset.
Training runs for 30 epochs with Adam~\cite{kingma2015adam} ($\text{lr}=10^{-4}$) and a
batch size of 256.
No environment interaction takes place; the model is evaluated directly after training.

\subsection{GAIL Training Loop}

GAIL is trained with the following configuration (Group~7 starter):

\begin{table}[h]
\centering
\caption{GAIL Hyperparameters}
\label{tab:hyperparams}
\begin{tabular}{lll}
\toprule
\textbf{Parameter} & \textbf{Value} & \textbf{Rationale} \\
\midrule
Total environment steps      & 5\,000\,000  & Budget parity with PPO \\
Rollout horizon $T$          & 1\,024       & Standard PPO Atari \\
PPO epochs per rollout       & 4            & Stable policy updates \\
PPO mini-batch size          & 128          & Memory constraint \\
Policy learning rate         & $2.5\times10^{-4}$ & Standard Atari PPO \\
Discriminator lr             & $3\times10^{-4}$   & Faster disc convergence \\
Disc.\ updates per rollout   & 3            & Avoids over-training disc \\
Discount $\gamma$            & 0.99         & Standard Atari \\
GAE $\lambda$                & 0.95         & Standard PPO \\
Clip $\epsilon$              & 0.2          & Standard PPO \\
Entropy coefficient $\lambda$& 0.01         & Residual exploration \\
Value loss coefficient       & 0.5          & Standard PPO \\
Max grad norm                & 0.5          & Gradient clipping \\
\bottomrule
\end{tabular}
\end{table}

At each rollout, the adversarial reward replaces the environment reward entirely.
The true environment reward is logged for monitoring but is never used in the
gradient computation.

\subsection{Advantage Estimation}

Generalised Advantage Estimation (GAE-$\lambda$)~\cite{schulman2016gae} is applied over
the adversarial rewards:

\begin{equation}
  \hat{A}_t = \sum_{l=0}^{T-t-1} (\gamma\lambda)^l \delta_{t+l},
  \qquad
  \delta_t = r_{\text{adv},t} + \gamma V(s_{t+1})(1-d_t) - V(s_t),
\end{equation}

where $d_t$ is the episode termination flag.

% ──────────────────────────────────────────────────────────────────────────────
\section{Experimental Protocol}
\label{sec:experiments}

\subsection{Budget Parity}

All algorithms are evaluated under a common step budget of 5\,000\,000 environment steps.
For GAIL, only agent interaction steps count; discriminator gradient steps are excluded.
This ensures a fair comparison of sample efficiency.

\subsection{Evaluation Protocol}

All algorithms are evaluated deterministically (greedy/argmax action selection) over
10 episodes with a fixed seed (\texttt{seed=0}).
The same environment wrappers and preprocessing pipeline are used for all three
algorithms.
Reported metrics are mean $\pm$ standard deviation of episode returns.

\subsection{Seeds}

\begin{itemize}
  \item Demonstration collection: seed 42.
  \item BC training: deterministic (depends only on data order).
  \item GAIL (primary runs): seeds 42, 0, 1.
  \item All evaluations: seed 0.
\end{itemize}

\subsection{Demonstration Ablation}

Two demonstration dataset sizes are tested as required:

\begin{table}[h]
\centering
\caption{Demonstration Datasets}
\label{tab:demos}
\begin{tabular}{llll}
\toprule
\textbf{Dataset} & \textbf{Steps} & \textbf{Episodes} & \textbf{Mean Return} \\
\midrule
\texttt{demos\_50k.npz} & 50\,000 & 296 & $78.1 \pm 16.2$ \\
\texttt{demos\_5k.npz}  &  5\,000 &  29 & $79.4 \pm 16.2$ \\
\bottomrule
\end{tabular}
\end{table}

The demonstrator policy is the best DQN checkpoint from Challenge~1, evaluated with
greedy (deterministic) action selection.
The mean return of the demonstrator during collection ($\approx78$--$79$ per game) reflects
the life-averaged scoring of the greedy policy under the evaluation wrappers.

% ──────────────────────────────────────────────────────────────────────────────
\section{Results}
\label{sec:results}

\subsection{Behavioral Cloning Training}

Table~\ref{tab:bc_training} summarises the BC training progression over 30 epochs on the
50\,000-step demonstration dataset (195 mini-batches per epoch, batch size 256).

\begin{table}[h]
\centering
\caption{BC Training — Selected Epoch Statistics}
\label{tab:bc_training}
\begin{tabular}{rrr}
\toprule
\textbf{Epoch} & \textbf{Cross-Entropy Loss} & \textbf{Accuracy} \\
\midrule
 1  & 1.9181 & 0.243 \\
 5  & 1.2557 & 0.543 \\
10  & 0.9919 & 0.638 \\
15  & 0.8027 & 0.720 \\
20  & 0.6552 & 0.776 \\
25  & 0.5340 & 0.820 \\
30  & 0.4460 & 0.852 \\
\bottomrule
\end{tabular}
\end{table}

The BC policy converges cleanly, achieving 85.2\% action-matching accuracy on the
demonstration data.
Training time is negligible compared to RL methods (no environment interaction required).

\subsection{Three-Way Performance Comparison}

Table~\ref{tab:comparison} presents the final evaluation results for all algorithms.
The evaluation follows the identical protocol: 10 deterministic episodes, seed 0,
same environment wrappers.

\begin{table}[h]
\centering
\caption{Final Evaluation Results — ALE/MsPacman-v5 (10 episodes, seed 0)}
\label{tab:comparison}
\begin{tabular}{lrrrr}
\toprule
\textbf{Algorithm} & \textbf{Mean} & \textbf{Std} & \textbf{Min} & \textbf{Max} \\
\midrule
BC (30 epochs, 50k demos)  & \textbf{866.0} & 280.2 &  490.0 & 1410.0 \\
DQN — Challenge 1          &  85.6          &  15.6 &   59.0 &  108.0 \\
GAIL (50k demos, seed 42)  &  70.0          &   0.0 &   70.0 &   70.0 \\
GAIL (5k demos, seed 42)   &  70.0          &   0.0 &   70.0 &   70.0 \\
\bottomrule
\end{tabular}
\end{table}

The results reveal three distinct outcome tiers:

\begin{enumerate}
  \item \textbf{BC} dominates all methods with a mean return of 866.0, demonstrating
        that deterministic imitation of the available demonstrations is highly effective
        for this game.
  \item \textbf{DQN} achieves a modest 85.6 mean return, consistent with the
        demonstrator's per-game performance (~78--85), indicating that the DQN policy
        has learned a basic but limited strategy.
  \item \textbf{GAIL} (both dataset sizes) collapses to exactly 70.0 with zero variance,
        corresponding to a degenerate fixed-action policy.
\end{enumerate}

\subsection{GAIL Collapse Analysis}

The GAIL score of $70.0 \pm 0.0$ across all 10 evaluation episodes indicates complete
policy collapse: the agent reproducibly reaches a terminal state that yields exactly 70
points regardless of episode stochasticity.
In Ms.\ Pac-Man, 70 points corresponds roughly to eating the first cluster of standard
pellets before dying, suggesting the policy learned a single local pattern and became
deterministically stuck.

The zero variance is the critical diagnostic indicator.
A healthy policy operating in a stochastic environment should exhibit episode-to-episode
variance; zero variance implies the policy has collapsed to a deterministic fixed point
that is invariant to the game's inherent stochasticity.

The root cause is \emph{discriminator saturation}: once the discriminator achieves near-
perfect separation between the expert (demonstrator) frames and the agent frames, it
outputs values close to 0 for all agent observations.
This drives $r_{\text{adv}} = \log D_\phi(s) \to -\infty$, which collapses the value
function and makes all PPO gradients uninformative.
Without a meaningful reward signal, the policy degrades to whatever local behavior
the initial random weights produce.

\subsection{Demonstration Size Ablation}

Both GAIL configurations (5k and 50k demonstrations) collapsed to identical performance
($70.0 \pm 0.0$), revealing that the failure mode is not caused by insufficient
demonstration data but rather by a fundamental instability in the adversarial training
dynamics on this visual, dense-reward environment.

% ──────────────────────────────────────────────────────────────────────────────
\section{Discussion}
\label{sec:discussion}

We address the four required analysis questions:

\subsection{Does GAIL outperform pure RL (DQN, PPO) on dense-reward games?}

No. In our experiment, GAIL significantly \emph{underperforms} both DQN (85.6) and BC
(866.0).
The key issue is that Ms.\ Pac-Man provides continuous, dense reward feedback.
DQN and PPO can exploit this signal immediately and directly, while GAIL replaces this
informative signal with a learned proxy (the adversarial reward) that is susceptible to
discriminator collapse.

This result is consistent with the theoretical analysis of GAIL: its advantages are
strongest in \emph{sparse-reward} settings where the environment reward is too infrequent
to drive policy gradient.
In dense-reward settings, replacing a reliable environment signal with a learned proxy
introduces unnecessary risk and complexity.
For games like Montezuma's Revenge or Private Eye (other groups' assignments), the
balance would likely shift in GAIL's favor.

\subsection{How sensitive is GAIL to demonstration quality and quantity?}

Our results show that GAIL is \emph{insensitive} to demonstration quantity in the
collapse regime: neither 5\,000 nor 50\,000 demonstrations prevented the adversarial
training from collapsing.
This insensitivity is itself informative — it indicates that the bottleneck is not
data coverage but rather the training dynamics of the discriminator.

The demonstration quality, assessed by the demonstrator's mean return ($\approx78$--$79$
per game), is modest but consistent.
The DQN demonstrator exhibits stable, reproducible behavior (std $\approx16$ across
datasets).
A higher-quality demonstrator (e.g., a trained PPO policy from Challenge~3) might produce
more informative demonstrations, but the discriminator collapse issue would persist unless
addressed architecturally.

\subsection{Does the adversarial reward remain informative, or does the discriminator collapse?}

The zero-variance GAIL evaluation results provide strong indirect evidence of discriminator
collapse.
In a healthy GAIL run, the discriminator should maintain accuracy close to 0.5
(perfect indistinguishability of agent from expert), providing a gradient signal that
continuously improves the policy.
Our discriminator appears to have converged to high accuracy (near-perfect separation),
making $r_{\text{adv}} = \log D_\phi(s) \approx \log\epsilon \approx -18$ for all agent
states, which is effectively a constant penalty with no gradient information.

Contributing factors include:
\begin{enumerate}
  \item \textbf{Visual domain complexity}: the $4\times84\times84$ pixel observations
        give the discriminator a very high-dimensional signal, making it easy for the
        discriminator to overfit early in training.
  \item \textbf{Dense-reward dynamics}: the expert demonstrations and the initial policy
        trajectories diverge rapidly because even random policies collect some pellets,
        but the demonstrator navigates with a clear strategy, creating a large early gap
        that the discriminator quickly learns to exploit.
  \item \textbf{Lack of gradient penalty}: we did not use gradient penalty regularization
        (suggested as optional in the challenge specification), which could have stabilized
        training by preventing the discriminator from saturating.
\end{enumerate}

\subsection{In what regime is BC competitive with full GAIL or PPO?}

Our results suggest BC is competitive with — and in this case, dominant over — full GAIL
when:
\begin{enumerate}
  \item \textbf{Dense rewards}: the game provides frequent feedback, making GAIL's proxy
        reward unnecessary.
  \item \textbf{High-quality demonstrations}: BC can only be as good as its demonstrations;
        here, the DQN demonstrator produces reliable movement patterns at $\approx79$
        per-game return.
  \item \textbf{Good demo coverage}: the primary dataset of 50\,000 steps from 296 episodes
        covers a diverse range of maze states, reducing distributional shift.
  \item \textbf{Limited compute}: BC trains entirely offline (no environment steps),
        making it orders of magnitude cheaper than GAIL or DQN for the same level of
        performance.
\end{enumerate}

The BC policy's high variance ($\pm280.2$) reveals the distributional shift effect in
practice: it achieves as high as 1410 and as low as 490, depending on whether the ghost
encounter situations fall within the demonstration distribution.
This variance would be reduced by additional supervised training epochs or by a DAgger-style
iterative dataset aggregation.

% ──────────────────────────────────────────────────────────────────────────────
\section{Limitations and Future Work}
\label{sec:limitations}

Several limitations constrain the current study:

\begin{enumerate}
  \item \textbf{Compute budget}: the 5M-step GAIL run was completed on a CPU
        (\texttt{CUDA available: False}), which significantly limited the ability to
        tune hyperparameters or add stabilization techniques such as gradient penalty.
  \item \textbf{PPO baseline absent}: Challenge~3 PPO results were not re-evaluated under
        the current framework, preventing a full four-way comparison.
        Future work should integrate the best PPO configuration from Challenge~3
        directly into this evaluation pipeline.
  \item \textbf{Single BC dataset size}: only the 50k demonstration dataset was used
        for BC; a BC vs.\ demo-size ablation (mirroring the GAIL ablation) would
        clarify how much of BC's success is due to data quantity vs.\ architecture.
  \item \textbf{No gradient penalty}: adding a gradient penalty term ($\lambda_{gp} = 10$)
        or spectral normalization to the discriminator~\cite{miyato2018sn} would likely
        prevent discriminator saturation in future runs.
  \item \textbf{State-action discriminator}: the \texttt{use\_action=True} variant of the
        discriminator was not tested; including action context may slow discriminator
        overfitting in image-based settings.
\end{enumerate}

% ──────────────────────────────────────────────────────────────────────────────
\section{Conclusion}
\label{sec:conclusion}

This paper presented a comparative study of DQN, Behavioral Cloning, and GAIL on the
\texttt{ALE/MsPacman-v5} environment.
The results confirm the Group~7 hypothesis in a particularly strong form: not only does
GAIL converge more slowly than pure RL methods on this dense-reward game, it fails to
converge at all due to discriminator collapse, while BC substantially outperforms both
GAIL and the DQN baseline.

The key lessons from this three-challenge series on Ms.\ Pac-Man are:
\begin{enumerate}
  \item \textbf{Dense rewards favor direct RL.} Both DQN and PPO benefit from the
        continuous game feedback that Ms.\ Pac-Man provides.
        GAIL's advantage is most pronounced in sparse-reward games where the true reward
        signal provides little gradient.
  \item \textbf{BC is a strong and underappreciated baseline.} When demonstrations
        are available and of reasonable quality, BC achieves competitive or superior
        performance with dramatically lower computational cost.
  \item \textbf{Discriminator stability is the central engineering challenge in GAIL.}
        Without explicit regularization (gradient penalty, spectral normalization,
        discriminator learning rate annealing), the discriminator easily saturates on
        complex visual domains, destroying the adversarial reward signal.
  \item \textbf{Algorithm choice should be guided by environment properties.} The
        three-challenge arc demonstrates that no single algorithm dominates across all
        settings: DQN excels with stable exploration, PPO provides on-policy stability,
        and GAIL is theoretically well-motivated for sparse-reward imitation but requires
        careful engineering to realize its potential.
\end{enumerate}

Future work will focus on GAIL stabilization through gradient penalty regularization,
the state-action discriminator variant, and integrating PPO results from Challenge~3
for a complete four-way algorithmic comparison.

% ──────────────────────────────────────────────────────────────────────────────
\begin{thebibliography}{00}

\bibitem{bellemare2013ale}
M.~G. Bellemare, Y.~Naddaf, J.~Veness, and M.~Bowling,
``The arcade learning environment: An evaluation platform for general agents,''
\textit{Journal of Artificial Intelligence Research}, vol. 47, pp. 253--279, 2013.

\bibitem{mnih2015dqn}
V.~Mnih, K.~Kavukcuoglu, D.~Silver, A.~A. Rusu, J.~Veness, M.~G. Bellemare,
A.~Graves, M.~Riedmiller, A.~K. Fidjeland, G.~Ostrovski, S.~Petersen, C.~Beattie,
A.~Sadik, I.~Antonoglou, H.~King, D.~Kumaran, D.~Wierstra, S.~Legg, and D.~Hassabis,
``Human-level control through deep reinforcement learning,''
\textit{Nature}, vol. 518, pp. 529--533, 2015.

\bibitem{schulman2017ppo}
J.~Schulman, F.~Wolski, P.~Dhariwal, A.~Radford, and O.~Klimov,
``Proximal policy optimization algorithms,''
\textit{arXiv preprint arXiv:1707.06347}, 2017.

\bibitem{ho2016gail}
J.~Ho and S.~Ermon,
``Generative adversarial imitation learning,''
in \textit{Advances in Neural Information Processing Systems (NeurIPS)}, 2016, pp. 4565--4573.

\bibitem{ross2011dagger}
S.~Ross, G.~Gordon, and D.~Bagnell,
``A reduction of imitation learning and structured prediction to no-regret online learning,''
in \textit{Proceedings of the 14th International Conference on Artificial Intelligence
and Statistics (AISTATS)}, 2011, pp. 627--635.

\bibitem{fu2018airl}
J.~Fu, K.~Luo, and S.~Levine,
``Learning robust rewards with adversarial inverse reinforcement learning,''
in \textit{International Conference on Learning Representations (ICLR)}, 2018.

\bibitem{schulman2016gae}
J.~Schulman, P.~Moritz, S.~Levine, M.~I. Jordan, and P.~Abbeel,
``High-dimensional continuous control using generalized advantage estimation,''
in \textit{International Conference on Learning Representations (ICLR)}, 2016.

\bibitem{goodfellow2014gan}
I.~Goodfellow, J.~Pouget-Abadie, M.~Mirza, B.~Xu, D.~Warde-Farley, S.~Ozair,
A.~Courville, and Y.~Bengio,
``Generative adversarial nets,''
in \textit{Advances in Neural Information Processing Systems (NeurIPS)}, 2014,
pp. 2672--2680.

\bibitem{kingma2015adam}
D.~P. Kingma and J.~Ba,
``Adam: A method for stochastic optimization,''
in \textit{International Conference on Learning Representations (ICLR)}, 2015.

\bibitem{torabi2018bco}
F.~Torabi, G.~Warnell, and P.~Stone,
``Behavioral cloning from observation,''
in \textit{Proceedings of the 27th International Joint Conference on Artificial
Intelligence (IJCAI)}, 2018, pp. 4950--4957.

\bibitem{miyato2018sn}
T.~Miyato, T.~Kataoka, M.~Koyama, and Y.~Yoshida,
``Spectral normalization for generative adversarial networks,''
in \textit{International Conference on Learning Representations (ICLR)}, 2018.

\end{thebibliography}

\end{document}
